%% 9-13 終章　工学の数学の地図（第15回：組 V9-1。ページ制限なし・偶数ページ。決定13）
%% ① 7本の路線と終点（series.py の ROUTES・ROUTE_STATIONS から写した）② 全9巻の路線図（\bkSeriesMap）③ 全部の駅
%% ④ 合流した5つの所（s=σ+jω、e^{At}、熱の行列の固有ベクトル＝sin、ガウスの形、波動方程式と音・光）
%% ⑤ 1つの現象を全巻の道具で（冷めるお湯・車のばね・棒の熱・ロボットの腕）⑥ 9-0 の回収（3つの部・くり返し出てくる場面）
%% ⑦ 名前と実際の操作の総まとめ ⑧ 専門書の読み方ガイド（記号の訳・RC の式を1つ読む）⑨ 物理シリーズへの橋 ⑩ ふりかえりと会話
%% 9-12 を指す所には「%% 9-12」の印をつけた（V9-6 の 9-12 ができたら、中身に合わせる）。
\makeatletter
\gdef\bk@partlabel{終章}\gdef\bk@partname{工学の数学の地図}
%% ===== kata9 の型（tools/kata9.md、組 V9-1）。頭文字 kixn は章ごとに置きかえる =====
% (1) 色：表計算のマスと、2つの世界の線
\colorlet{kixnN}{bkOrange!16}% 1行め（成分の番号・周波数）のマス
\colorlet{kixnK}{bkGreen!16}% 2行め（成分の大きさ）のマス
\colorlet{kixnS}{bkYellow!45}% 合計・注目するマス
\colorlet{kixnW}{bkBlue!80!black}% 時間の世界の線（波形・動き）
\colorlet{kixnF}{bkOrange!85!black}% 周波数の世界の棒
\colorlet{kixnP}{bkRed!80!black}% s の世界の点（×）
% (2) 表計算の数式バー  \kixnfx{セル}{式}（式の中の $ は \$、^ は \^{}）
\newcommand\kixnfx[2]{\tikz[baseline=(f.base)]{\node[draw=bkLine, line width=0.6pt, fill=white, minimum height=5.5mm, inner xsep=2mm, font=\sffamily\footnotesize, minimum width=12mm] (c) {#1};
  \node[anchor=west, draw=bkLine, line width=0.6pt, fill=white, minimum height=5.5mm, inner xsep=2mm, minimum width=78mm, text width=74mm, align=left, font=\ttfamily\footnotesize] (f) at ([xshift=1mm]c.east) {#2};}}
% (3) グラフのスタイル（どれも scale only axis：はばと高さは軸の箱の大きさ）
\pgfplotsset{
  kixn base/.style={bkplot, scale only axis, width=47mm, height=28mm, axis lines=left, grid=none,
    every axis x label/.style={at={(axis description cs:0.5,0)}, anchor=north, yshift=-5mm, font=\sffamily\small, text=bkInk},
    every axis y label/.style={at={(axis description cs:0,1)}, anchor=south, yshift=1.5mm, font=\sffamily\small, text=bkInk}},
  kixn wave/.style={kixn base, extra y ticks={0}, extra y tick labels={0}, extra y tick style={grid=major, major grid style={bkInk!45, line width=0.5pt}}},
  kixn spec/.style={kixn base, ymin=0},
  kixn splane/.style={bkplot, scale only axis, width=47mm, height=47mm, axis equal image, grid=major,
    every axis x label/.style={at={(ticklabel* cs:1)}, anchor=north east, yshift=-0.5mm, font=\sffamily\small, text=bkInk},
    every axis y label/.style={at={(ticklabel* cs:1)}, anchor=south, yshift=1mm, font=\sffamily\small, text=bkInk}},
  kixn sig/.style={kixnW, line width=1.4pt, no markers, samples=241, smooth},
  kixn part/.style={bkGreen!45!black, line width=0.8pt, dash pattern=on 3pt off 1.5pt, no markers, samples=241, smooth},
  kixn bar/.style={ycomb, draw=kixnF, line width=3pt, no markers},
  kixn dot/.style={only marks, mark=*, mark size=1.3pt, kixnW}}
% (4) 2つの世界の見出しと、行き来の矢印  \kixnworlds[見出しの高さ]{左の題}{右の題}{上の矢印（→）}{下の矢印（←）}
%     左の軸は at={(0mm,0mm)}、右の軸は at={(76mm,0mm)}（はば 47 mm）に置く。見出しは軸の上（ふつう 42 mm）
\tikzset{
  kixn world/.style={draw, line width=1pt, rounded corners=1.5mm, fill=white, font=\sffamily\bfseries\small, text=bkInk,
    minimum width=30mm, minimum height=7mm, inner xsep=2mm},
  kixn bridge/.style={-{Stealth[length=2.4mm]}, line width=1pt, bkInk!80},
  kixn blab/.style={font=\sffamily\footnotesize, text=bkInk, inner sep=0.8mm}}
\newcommand\kixnworlds[5][42mm]{%
  \node[kixn world, draw=kixnW] (kixnL) at (23.5mm,#1) {#2};
  \node[kixn world, draw=kixnF] (kixnR) at (99.5mm,#1) {#3};
  \draw[kixn bridge] ([yshift=1.3mm, xshift=1.5mm]kixnL.east) -- node[kixn blab, above] {#4} ([yshift=1.3mm, xshift=-1.5mm]kixnR.west);
  \draw[kixn bridge] ([yshift=-1.3mm, xshift=-1.5mm]kixnR.west) -- node[kixn blab, below] {#5} ([yshift=-1.3mm, xshift=1.5mm]kixnL.east);}
% (5) s の世界の点（axis の中で）  \kixnpole{実部}{虚部}（×）  \kixnzero{実部}{虚部}（○）
\newcommand\kixnpole[3][kixnP]{\draw[#1, line width=1.3pt] (axis cs:#2,#3) ++(-1.3mm,-1.3mm) -- ++(2.6mm,2.6mm) (axis cs:#2,#3) ++(-1.3mm,1.3mm) -- ++(2.6mm,-2.6mm);}
\newcommand\kixnzero[3][kixnP]{\draw[#1, line width=1.1pt, fill=white] (axis cs:#2,#3) circle[radius=1.2mm];}
%% ===== この章だけで使う部品（頭文字 kixn） =====
\tikzset{
  kixn lab/.style={font=\sffamily\footnotesize, inner sep=1pt, text=bkInk},
  kixn arr/.style={-{Stealth[length=2.2mm, width=1.9mm]}, line width=1pt, bkInk!75},
  kixn eq/.style={draw=\bkTT, line width=0.9pt, rounded corners=1.5mm, fill=bkTheme!6, inner xsep=2mm, inner ysep=1.8mm, align=flush center, font=\small, text=bkInk},
  kixn st/.style={draw=bkInk!45, line width=0.8pt, rounded corners=1.2mm, fill=white, inner xsep=1.6mm, inner ysep=1.4mm,
    minimum width=29.2mm, anchor=north, minimum height=25mm},
  kixn stt/.style={text width=26mm, align=flush left, font=\footnotesize, text=bkInk, anchor=north west, inner sep=0pt}}
% 合流の図の左の1行  \kixnin{路線番号}{章}{ことば}
\newcommand\kixnin[3]{\setbox0\hbox{\bkrouteicon[1.2em]{#1}\enspace{\sffamily\bfseries #2}\enspace}\hangindent=\wd0\hangafter=1\noindent\box0 #3\par\vspace{1.5pt}}
% 合流の図  \kixnmerge{左（\kixnin の行）}{合流した式}{合流した駅}
\newcommand\kixnmerge[3]{%
  \begin{tikzpicture}[x=1mm, y=1mm]
    \node[anchor=west, text width=64mm, align=left, font=\small, inner sep=0pt] (L) at (0,0) {#1};
    \node[kixn eq, anchor=west, text width=50mm] (R) at (75,0) {#2};
    \draw[kixn arr] ([xshift=1.5mm]L.east |- R.west) -- ([xshift=-1.5mm]R.west);
    \node[anchor=south, font=\sffamily\bfseries\footnotesize, text=\bkTT, inner sep=0.8mm] at (R.north) {合流した駅：#3};
  \end{tikzpicture}}
% 1つの現象の道のり（4つずつ折り返す）  \kixnchain{{章}/{ことば}, {章}/{ことば}, …}{いくつめまで}
\newcommand\kixnchain[2]{%
  \begin{tikzpicture}[x=1mm, y=1mm]
    \foreach \kixnc/\kixnt [count=\kixni from 0] in {#1} {
      \pgfmathtruncatemacro\kixnr{\kixni/4}\pgfmathtruncatemacro\kixnk{mod(\kixni,4)}
      \node[kixn st] (kixnb\kixni) at ({\kixnk*33.6+14.6},{-\kixnr*33}) {};
      \node[kixn stt] at ([xshift=1.6mm, yshift=-1.6mm]kixnb\kixni.north west) {{\sffamily\bfseries\color{\bkTT}\kixnc}\par\vspace{1pt}\kixnt};}
    \foreach \kixni in {1,...,#2} {
      \pgfmathtruncatemacro\kixnp{\kixni-1}\pgfmathtruncatemacro\kixnk{mod(\kixni,4)}
      \ifnum\kixnk=0
        \draw[kixn arr, rounded corners=2mm] ([yshift=1.5mm]kixnb\kixnp.south) ++(0,-1.5) -- ++(0,-2.6) -| ([yshift=1.5mm]kixnb\kixni.north);
      \else
        \draw[kixn arr] ([xshift=0.8mm]kixnb\kixnp.east |- kixnb\kixni.north) ++(0,-12.5) -- ++(4.2,0);
      \fi}
  \end{tikzpicture}}
% 箇条書きの1項目  \kixnitem{文}
\newcommand\kixnitem[1]{\par\noindent\hangindent=1\zw\hangafter=1 ・#1\par}
% 会話の1行  \kixntalk{y}{sora|mina|delta}{名前}{せりふ}
\newcommand\kixntalk[4]{%
  \bkChar{#2}{(7,#1-6.5)}{12mm}
  \node[rectangle callout, callout absolute pointer={(15.5,#1)}, anchor=west, draw=bkTheme!80!black, line width=0.8pt,
    fill=white, rounded corners=2.5mm, inner xsep=3mm, inner ysep=1.8mm, text width=103mm, align=left,
    font=\small] at (21,#1) {#4};
  \node[anchor=north, font=\sffamily\footnotesize, text=bkGray, inner sep=0pt] at (7,#1-7.5) {#3};}
% 7本の路線の星（はてな？の絵）  \kixnstar{半径}
\newcommand\kixnstar[1]{%
  \foreach \kixni/\kixna in {1/90,2/141.4,3/192.9,4/244.3,5/295.7,6/347.1,7/38.6} {\draw[bkRoute\kixni, line width=2pt, line cap=round] (\kixna:#1) -- (0,0);
    \node[inner sep=0pt] at (\kixna:#1+4.2) {\bkrouteicon[5.4mm]{\kixni}};}
  \filldraw[fill=white, draw=bkInk, line width=0.9pt] (0,0) circle (3);
  \node[font=\sffamily\bfseries\footnotesize, text=bkInk] at (0,0) {?};}
% 見出しのアイコン（路線：線と2つの駅）
\newcommand\bk@icon@kixnroute{%
  \draw[bkOn, line width=0.13*\bk@iu, line cap=round] (-0.62,0) -- (0.62,0);
  \filldraw[fill=white, draw=bkOn, line width=0.09*\bk@iu] (-0.36,0) circle (0.2) (0.36,0) circle (0.2);}
\makeatother
\makeatletter
%% ===== 「名前と操作」の型（tools/kata9-int.md、組 V9-2）。頭文字 kixn は章ごとに置きかえる =====
\colorlet{kixnNc}{bkOrange!85!black}% 名前の枠の線（点線）
\colorlet{kixnSc}{bkTheme}% 実際の操作の枠の線（実線）
\colorlet{kixnSb}{bkTheme!8}% 実際の操作の地
\tikzset{
  kixn dan/.style={font=\sffamily\bfseries\footnotesize, rounded corners=1.2mm, inner xsep=1.6mm, inner ysep=0.5mm, line width=0.9pt},
  kixn dan1/.style={kixn dan, draw=kixnNc, dash pattern=on 2.2pt off 1.4pt, fill=white, text=bkInk},
  kixn dan2/.style={kixn dan, draw=kixnSc, fill=kixnSb, text=bkInk},
  kixn dan3/.style={kixn dan, draw=bkTheme, fill=bkTheme, text=bkOn},
  kixn name/.style={draw=kixnNc, line width=0.9pt, dash pattern=on 2.2pt off 1.4pt, rounded corners=1.2mm, inner xsep=1.2mm, inner ysep=1.2mm, fill=white},
  kixn seki/.style={draw=kixnSc, line width=0.9pt, rounded corners=1.2mm, inner xsep=1.2mm, inner ysep=1.2mm, fill=kixnSb},
  kixn tag/.style={font=\sffamily\bfseries\footnotesize, text=bkInk, inner sep=0.4mm},
  kixn q/.style={diamond, aspect=4.6, draw=bkYellow!75!black, line width=0.9pt, fill=bkYellowBg, font=\sffamily\footnotesize, align=flush center, inner sep=0.2mm, text=bkInk},
  kixn do/.style={rectangle, draw=\bkTT, line width=0.9pt, fill=white, rounded corners=0.6mm, font=\sffamily\footnotesize, align=flush center, inner xsep=1.5mm, inner ysep=1.2mm, text=bkInk},
  kixn go/.style={-{Stealth[length=2.2mm]}, line width=0.9pt, bkInk!75, shorten <=1.5mm, shorten >=1.5mm},
  kixn yn/.style={font=\sffamily\bfseries\footnotesize, text=bkInk, fill=white, inner sep=0.5mm},
}
% 実際の操作の4種類の名前（実線の札の字）  1 ふつうの定積分　2 表を引く　3 部分分数　4 行列の計算
\def\kixn@op#1{\ifcase#1\or ふつうの定積分\or 表を引く\or 部分分数\or 行列の計算\fi}
% 段の札（例題の中で）  \kixndan{1|2|3}
\newcommand\kixndan[1]{\tikz[baseline=(d.base)]\node[kixn dan#1] (d) {\vphantom{あ}\ifcase#1\or 名前\or 直す\or 楽をする\fi};}
% 直す段の札に、どの道具かをそえる（例題の中で）  \kixnnaosu{1|2|3|4} →「直す：表を引く」
\newcommand\kixnnaosu[1]{\tikz[baseline=(d.base)]\node[kixn dan2] (d) {\vphantom{あ}直す：\kixn@op{#1}};}
% 実線の札だけ（表の「手順」の頭・本文）  \kixnop{1|2|3|4}
\newcommand\kixnop[1]{\tikz[baseline=(d.base)]\node[kixn dan2] (d) {\vphantom{あ}\kixn@op{#1}};}
% 段の見出し（本文で）  \kixndanhead{1|2|3}{見出し}
\newcommand\kixndanhead[2]{\par\addvspace{8pt}\Needspace{7\baselineskip}\noindent
  \tikz[baseline=(d.base)]\node[kixn dan#1, font=\sffamily\bfseries\small, inner ysep=0.8mm] (d) {\vphantom{あ}\ifcase#1\or ①\,名前の段\or ②\,直す段\or ③\,楽をする段\fi};%
  \hspace{0.7\zw}{\sffamily\bfseries\color{\bkTT}#2}\par\nobreak\vspace{3pt}\noindent\ignorespaces}
% 場合の札  \kixncase{C}（′ は \kixncase{A$'$}）
\newcommand\kixncase[1]{\tikz[baseline=(c.base)]\node[fill=bkTheme, text=bkOn, font=\sffamily\bfseries\small, rounded corners=0.8mm, inner xsep=1mm, inner ysep=0.6mm, minimum width=5.4mm] (c) {\vphantom{Ag}#1};}
% 例題の1行め  \kixnbaai{ことば}{9-2}{C}（場合が2つなら {D,B}：通った順）
\newcommand\kixnbaai[3]{{\sffamily\bfseries 場合：}#1（表#2の\kixn@cases#3,\@nil）}
\def\kixn@cases#1,#2\@nil{\kixncase{#1}\ifx\relax#2\relax\else と\kixn@cases#2\@nil\fi}
% 式の中の「名前」（点線）と「実際の操作」（実線）
\newcommand\kixnN[1]{\tikz[baseline=(n.base)]\node[kixn name] (n) {$\displaystyle #1$};}
\newcommand\kixnNt[1]{\tikz[baseline=(n.base)]{\node[kixn name] (n) {$\displaystyle #1$};\node[kixn tag, anchor=north] at ([yshift=-0.8mm]n.south) {名前};}}
\newcommand\kixnS[1]{\tikz[baseline=(n.base)]\node[kixn seki] (n) {$\displaystyle #1$};}
% 札つきの実線の枠  \kixnSt{1|2|3|4}{式}（札の字：ふつうの定積分・表を引く・部分分数・行列の計算）
\newcommand\kixnSt[2]{\tikz[baseline=(n.base)]{\node[kixn seki] (n) {$\displaystyle #2$};\node[kixn tag, anchor=north] at ([yshift=-0.8mm]n.south) {\kixn@op{#1}};}}
% 場合分けの表の題の帯  \kixncasetitle{9-2}{題}
\newcommand\kixncasetitle[2]{\par\addvspace{6pt}\Needspace{8\baselineskip}\noindent\begin{tikzpicture}[baseline=(t.base)]
  \node[fill=bkTheme, text=bkOn, font=\sffamily\bfseries\small, rounded corners=1mm, inner xsep=2mm, inner ysep=1mm] (t) {場合分けの表\,#1};
  \node[anchor=west, font=\sffamily\bfseries, text=\bkTT, inner sep=0pt] at ([xshift=2mm]t.east) {#2};
  \end{tikzpicture}\par\nobreak\vspace{3pt}}
% 名前と操作の対応表の小さな絵  \kixnpic{1〜7}
%  1 フーリエ係数（1周期で cos をかける）　2 複素係数（回る矢印）　3 フーリエ変換（−∞〜∞）
%  4 ラプラス変換（0から、e^{-st} で小さく）　5 逆ラプラス変換（s の平面の極）　6 伝達関数（箱）　7 e^{At}（行列とうずまき）
\newcommand\kixnpic[1]{\begin{tikzpicture}[x=1mm, y=1mm, baseline=(current bounding box.center)]
  \ifcase#1\or
  \draw[bkGray, line width=0.4pt] (0,4) -- (14,4);
  \draw[bkInk, line width=0.8pt] (1,4) -- (1,7) -- (7,7) -- (7,1) -- (13,1) -- (13,4);
  \draw[bkTheme, line width=0.8pt, dash pattern=on 1.6pt off 0.9pt, domain=1:13, samples=40, smooth] plot (\x,{4+2.6*cos(30*(\x-1))});
  \draw[bkInk, line width=0.5pt] (1,-0.6) -- (1,-1.6) -- (13,-1.6) -- (13,-0.6);
  \or
  \draw[bkGray, line width=0.4pt] (2,4) -- (12,4) (7,-0.5) -- (7,8.5);
  \draw[bkInk, line width=0.7pt] (7,4) circle (3.6);
  \draw[bkTheme, line width=1.1pt, -{Stealth[length=1.6mm]}] (7,4) -- ({7+3.6*cos(-40)},{4+3.6*sin(-40)});
  \draw[bkInk, line width=0.6pt, -{Stealth[length=1.3mm]}] ({7+5*cos(10)},{4+5*sin(10)}) arc (10:-60:5);
  \or
  \draw[bkInk, line width=0.6pt, {Stealth[length=1.4mm]}-{Stealth[length=1.4mm]}] (0,1) -- (14,1);
  \draw[bkInk, line width=0.8pt] (1.5,1) -- (5.5,1) -- (5.5,6.5) -- (8.5,6.5) -- (8.5,1);
  \draw[bkTheme, line width=0.8pt, dash pattern=on 1.6pt off 0.9pt, domain=1.5:12.5, samples=50, smooth] plot (\x,{1+2.2*(1+cos(80*(\x-7)))/2});
  \or
  \draw[bkInk, line width=0.6pt, -{Stealth[length=1.4mm]}] (1,1) -- (14,1);
  \draw[bkInk, line width=0.6pt] (1,0) -- (1,8.5);
  \fill[bkTheme!25, domain=1:13, samples=40, smooth] (1,1) -- plot (\x,{1+6.5*exp(-0.3*(\x-1))}) -- (13,1) -- cycle;
  \draw[bkTheme, line width=0.8pt, domain=1:13, samples=40, smooth] plot (\x,{1+6.5*exp(-0.3*(\x-1))});
  \or
  \draw[bkGray, line width=0.4pt] (0,4) -- (14,4) (10,-0.5) -- (10,8.5);
  \foreach \kixnx/\kixny in {3/4,6.5/6.8,6.5/1.2} {\draw[bkRed, line width=1pt] ({\kixnx-1},{\kixny-1}) -- ({\kixnx+1},{\kixny+1}) ({\kixnx-1},{\kixny+1}) -- ({\kixnx+1},{\kixny-1});}
  \or
  \draw[bkInk, line width=0.6pt, -{Stealth[length=1.4mm]}] (0,4) -- (4,4);
  \draw[bkTheme, line width=0.8pt, fill=bkTheme!12] (4.4,1) rectangle (9.6,7);
  \node[font=\footnotesize, text=bkInk, inner sep=0pt] at (7,4) {$G$};
  \draw[bkInk, line width=0.6pt, -{Stealth[length=1.4mm]}] (10,4) -- (14,4);
  \or
  \draw[bkInk, line width=0.6pt] (1.6,0.5) -- (0.6,0.5) -- (0.6,7.5) -- (1.6,7.5) (5.4,0.5) -- (6.4,0.5) -- (6.4,7.5) -- (5.4,7.5);
  \foreach \kixnx/\kixny in {2.3/2.3,4.7/2.3,2.3/5.7,4.7/5.7} {\fill[bkInk] (\kixnx,\kixny) circle (0.5);}
  \draw[bkTheme, line width=0.8pt, domain=0:600, samples=80, smooth] plot ({10.5+3.4*exp(-0.0035*\x)*cos(\x)},{4+3.4*exp(-0.0035*\x)*sin(\x)});
  \fi
\end{tikzpicture}}
\makeatother

\bkchapter{9}{13}{工学の数学の地図}{\bklevelrange{3}{5}}{工学の数学の地図}

\begin{goal}
\begin{itemize}[label={\color{\bkTT}●}, leftmargin=1.6\zw, itemsep=0pt, topsep=1pt]
\item \bklevel{3}\enspace 全9巻の7本の路線が、どの駅を通って、どこを終点にしたかを、地図でたどれる。
\item \bklevel{4}\enspace 路線が合流した5つの所で、どの巻の道具とどの巻の道具が出会ったかを、式と例で説明できる。
\item \bklevel{5}\enspace 専門書の式を、名前の段で読み、記号を訳し、4つの道具のどれで計算するかを決められる。
\end{itemize}
\end{goal}


\begin{hatena}
\noindent\begin{minipage}[c]{0.58\linewidth}
\raggedright
このシリーズの章は、どれも7本の路線のどれかの駅でした。7本の路線は、第1巻・第2巻の別々の駅から出発しました。\par\vspace{4pt}
その7本は、第9巻までの間に、どこかで出会ったのでしょうか。どこで、どの路線とどの路線が合流したのかな？
\end{minipage}\hfill
\begin{minipage}[c]{0.38\linewidth}
\centering
\begin{tikzpicture}[x=1mm, y=1mm]
  \kixnstar{17}
\end{tikzpicture}
\end{minipage}
\end{hatena}

小学1年の「リンゴは何個？」から始まったこのシリーズも、この章が最後です。全9巻、152の章を読んできました。数を数えることから始めて、かけ算、グラフ、三角関数、微分と積分、微分方程式、行列、ベクトル場、そして第9巻のフーリエ解析とラプラス変換まで、一段ずつ積み上げてきました。ここまで来たあなたは、大学の工学部で使う数学の、大事な道具をひととおり手に入れています。本当によくがんばりました。

この章では、問題は解きません。これまでに乗ってきた路線を、地図にしてふり返ります。

この章は、どこから読んでもかまいません。前の巻の章の番号がたくさん出てきます。なつかしい章があれば、ぜひ開いてみてください。1回目に読んだときより、ずっとよくわかるはずです。

\begin{omoidasou}
\bkrecall{9-0}{第9巻は、3つの部と終章。路線④・⑦・③・⑤の4本が、この巻に終点をもつ。}
\bkrecall{9-0}{記号は、まだ計算していない「名前」。実際の値は、4つの道具（ふつうの定積分・表を引く・部分分数・行列の計算）のどれかに直して出す。}
\bkrecall{8-0}{第8巻では、路線①・②・⑥の3本が終点に着いた（8-4・8-11・8-5）。}
\bkrecall{6-0}{第6巻の微分方程式は、路線①〜④が乗りかえる駅（全9巻の路線図の点線の囲み）。}
\end{omoidasou}

%% ================= 7本の路線 =================
\bksection{kixnroute}{7本の路線をふり返る}

\bksubsection{1}{7本の路線と、その終点}

\label{bk:9-13:rosen}このシリーズの章は、どれも7本の路線のどれかに乗っています。路線は「1つの考え方の道」です。同じ路線の駅は、小学校の算数から大学の数学まで、1つの考え方でつながっています。どの路線も、最後は工学で使う分野（終点）に着きます。

\begin{bkcenter}
\begin{bktable}{colspec={Q[c,m,wd=7mm] Q[l,m,wd=20mm] Q[l,m,wd=66mm] Q[l,m,wd=26mm]}, rowsep=1.4pt, colsep=3pt, cells={font=\small}, row{1}={font=\sffamily\bfseries\small, bg=bkTheme!10}}
 & 路線 & 大きな駅のつながり & 終点 \\
\bkrouteicon{1} & 変化 & ひき算（1-5）→ 四則（2-3）→ グラフの傾き（3-5）→ 瞬間の速さ（5-1）→ 微分方程式（6-1）→ \textbf{勾配 grad}（8-4） & 流体・熱 \\
\bkrouteicon{2} & たし合わせ & 面積（1-21）→ 四則（2-3）→ 細かく分けてたす（4-5）→ 細い長方形（5-10）→ 積分で解く（6-3）→ \textbf{ガウスの定理}（8-11） & 電磁気 \\
\bkrouteicon{3} & 倍々 & 累乗と指数（2-7）→ 指数関数（3-11）→ $e$（5-14）→ 冷めるお湯（6-4）→ \textbf{ラプラス変換}（9-5） & 制御 \\
\bkrouteicon{4} & 回転と波 & 角の大きさ（1-20）→ 三角関数（4-8）→ $\sin$ の微分（5-15）→ オイラーの公式（6-10）→ \textbf{フーリエ級数}（9-1） & 音・交流 \\
\bkrouteicon{5} & 表 & 九九の表（1-13）→ 表計算（2-18）→ 連立方程式（3-6）→ 行列（7-1）→ \textbf{状態方程式}（9-10） & 回路網・制御 \\
\bkrouteicon{6} & 矢印 & 数直線（2-2）→ 座標（3-2）→ ベクトル（4-10）→ 3次元の回転（7-7）→ \textbf{ベクトル場}（8-5） & ロボット・電磁気 \\
\bkrouteicon{7} & はかる & 目もり（1-11）→ ずれ（2-6）→ 標準偏差（3-14）→ \textbf{フーリエ変換}（9-4） & 計測・診断 \\
\end{bktable}
\bktabcaption{7本の路線の大きな駅と終点（太字が終点の駅）。2-3「＋－×÷は、何を表しているの？」は、①と②の両方の駅。}
\label{bk:9-13:tab-rosen}
\end{bkcenter}

7本の路線の名前は、それぞれの考え方の道を表しています。
\begin{multicols}{2}\setlength{\columnsep}{6mm}
\noindent\bkrouteicon{1}\enspace\textbf{変化}：ひき算→傾き→微分\par
\noindent\bkrouteicon{2}\enspace\textbf{たし合わせ}：かけ算→面積→積分\par
\noindent\bkrouteicon{3}\enspace\textbf{倍々}：累乗→指数→$e$\par
\noindent\bkrouteicon{4}\enspace\textbf{回転と波}：角→$\sin$→複素数\par
\noindent\bkrouteicon{5}\enspace\textbf{表}：九九の表→行列\par
\noindent\bkrouteicon{6}\enspace\textbf{矢印}：向き→ベクトル→場\par
\noindent\bkrouteicon{7}\enspace\textbf{はかる}：目もり→誤差→統計\par
\end{multicols}

\bksubsection{2}{全9巻の路線図}

\label{bk:9-13:rosenzu}どの巻の初めにもあった「全9巻の路線図」を、もう一度見てみましょう。下へ行くほど先の巻です。いまは、いちばん下の第9巻にいます。

\begin{bkfigure}
\resizebox{112mm}{!}{\makeatletter\bkSeriesMap\makeatother}
\bkfigcaption{全9巻の路線図（どの巻の初めにもある図と同じ）。たての線が7本の路線、横の行が巻、丸が主な駅。いちばん下が終点。第9巻の行（色の帯）では、③・④・⑤・⑦の4本の路線が終点に着く。}
\label{bk:9-13:fig-rosenzu}
\end{bkfigure}

路線図のたての線をたどると、どの路線も第1巻か第2巻（小学1年〜中学1年）から出発して、第8巻か第9巻で終点に着いています。第6巻の点線の囲みは「乗りかえ駅」でした。①〜④の4本の路線が、第6巻の微分方程式でいっせいに出会ったのです（6-0）。

\bksubsection{3}{全部の駅}

\label{bk:9-13:eki}路線図にかいた駅は、主な駅だけです。全部の駅（章）を並べると、次のようになります。数字は章の番号で、左から順に、前の巻から先の巻へ並んでいます（終章の 9-13 は7本すべての駅なので、この表には入れていません。2-3 は①と②の両方に入っています）。

\begin{bkcenter}
\begin{bktable}{colspec={Q[c,m,wd=7mm] Q[l,m,wd=17mm] X[l,m] Q[c,m,wd=9mm]}, rowsep=1.4pt, colsep=3pt, cells={font=\footnotesize}, row{1}={font=\sffamily\bfseries\small, bg=bkTheme!10}}
 & 路線 & 駅（章の番号） & 駅の数 \\
\bkrouteicon{1} & 変化 & 1-5・1-10・1-14・1-18・2-0・2-3・2-4・2-15・2-16・2-19・3-5・3-7・3-8・3-9・3-10・5-0・5-1・5-2・5-3・5-4・5-5・5-6・5-7・5-8・5-17・5-18・6-0・6-1・6-2・6-5・6-15・8-1・8-2・8-3・8-4・8-12・8-13 & 37 \\
\bkrouteicon{2} & たし合わせ & 1-1・1-4・1-7・1-12・1-21・2-3・4-5・5-9・5-10・5-11・5-12・5-13・5-16・6-3・8-6・8-7・8-8・8-11 & 18 \\
\bkrouteicon{3} & 倍々 & 1-2・2-7・2-8・2-17・3-11・3-12・5-14・6-4・6-11・9-5・9-6・9-7・9-8・9-9 & 14 \\
\bkrouteicon{4} & 回転と波 & 1-15・1-19・1-20・4-0・4-2・4-6・4-7・4-8・4-9・5-15・6-6・6-7・6-8・6-9・6-10・6-12・6-13・7-6・9-0・9-1・9-2・9-3 & 22 \\
\bkrouteicon{5} & 表 & 1-13・1-17・2-11・2-12・2-13・2-14・2-18・3-3・3-6・6-14・7-0・7-1・7-2・7-3・7-4・7-5・7-8・7-9・7-10・7-11・9-10・9-11 & 22 \\
\bkrouteicon{6} & 矢印 & 1-3・1-6・2-2・2-9・3-2・4-1・4-3・4-4・4-10・4-11・4-12・4-13・4-14・4-15・7-7・8-0・8-5・8-9・8-10・8-14 & 20 \\
\bkrouteicon{7} & はかる & 1-0・1-8・1-9・1-11・1-16・2-1・2-5・2-6・2-10・3-0・3-1・3-4・3-13・3-14・3-15・3-16・3-17・9-4・9-12 & 19 \\
\end{bktable}
\bktabcaption{7本の路線の全部の駅。いちばん駅が多いのは①変化（37駅）。微分は、算数から大学まで、いちばん長く育ってきた考え方。}
\end{bkcenter}

いちばん長い路線は①変化です。「きのうより、どれだけのびた？」（1-5 のひき算）が、「グラフの傾き」（3-5）になり、「瞬間の速さ」（5-1）になり、「いちばん急な向き」（8-4）になりました。1つの考え方が、学年が上がるたびに、より広い場面で使えるように育っていったのです。

\begin{bkcenter}
\begin{bktable}{colspec={Q[c,m,wd=14mm] Q[l,m,wd=40mm] Q[l,m,wd=52mm] Q[c,m,wd=13mm]}, rowsep=1.2pt, colsep=3pt, cells={font=\small}, row{1}={font=\sffamily\bfseries\small, bg=bkTheme!10}}
巻 & 題名 & 副題 & 章の数 \\
第1巻 & かぞえる・はかる & & 22 \\
第2巻 & 数と量のことば & & 20 \\
第3巻 & 関数とデータ & きまりを見つける & 18 \\
第4巻 & 形と向き & 図形・三角関数・ベクトル & 16 \\
第5巻 & 変化をつかむ & 微分と積分 & 19 \\
第6巻 & 未来を計算する & 微分方程式と複素数 & 16 \\
第7巻 & まとめて計算する & 行列 & 12 \\
第8巻 & 空間に広がる量 & ベクトル解析 & 15 \\
第9巻 & 波と変化を読みとく & フーリエ解析とラプラス変換 & 14 \\
\end{bktable}
\bktabcaption{全9巻の題名と章の数（序章・終章をふくむ）。あわせて152章。}
\end{bkcenter}

%% ================= 合流 =================
\bksection{doguu}{路線が合流した所}

\bksubsection{4}{合流とは}

\label{bk:9-13:goryu}全9巻の路線図では、7本の路線は、たてにまっすぐ並んでいました。けれども、本当は、路線どうしは何回も出会っています。ある路線で作った道具と、別の路線で作った道具が、1つの式の中で組み合わさる所があるのです。この本では、それを路線の\textbf{合流}（ごうりゅう）とよびます。\bkidx{ごうりゅう}{合流（路線の）}

\noindent\textbf{\bkhatenaref の答え}\par
大きな合流は、次の5つです。どれも、第6巻から第9巻の間で起きました。

\begin{bkcenter}
\begin{bktable}{colspec={Q[c,m,wd=8mm] Q[l,m,wd=40mm] Q[l,m,wd=44mm] Q[l,m,wd=30mm]}, rowsep=1.4pt, colsep=3pt, cells={font=\small}, row{1}={font=\sffamily\bfseries\small, bg=bkTheme!10}}
 & 合流した所 & 出会った路線 & 合流した駅 \\
（ア） & $s=\sigma+j\omega$ & ③倍々と④回転と波 & 6-12・9-5・9-9 \\
（イ） & 行列の指数関数 $e^{At}$ & ③倍々と⑤表 & 7-10・9-11 \\
（ウ） & 熱の行列の固有ベクトルは $\sin$ の形 & ①変化と④回転と波と⑤表 & 9-11 \\
（エ） & ガウスの形 & ②たし合わせと④回転と波と⑦はかる & 9-4 \\
（オ） & 波動方程式と、音・光 & ①変化と⑥矢印と④回転と波 & 9-0・9-1 \\
\end{bktable}
\bktabcaption{路線が合流した5つの所（くわしくは⑤〜⑨）}
\label{bk:9-13:tab-goryu}
\end{bkcenter}

表\ref{bk:9-13:tab-goryu}の合流を、路線の地図にかき入れると、図\ref{bk:9-13:fig-goryu}のようになります。どの合流も、第6巻から先の巻で起きています。

\begin{bkfigure}
\begin{tikzpicture}[x=1mm, y=1mm]
  % 巻の列
  \foreach \kixnv/\kixnx/\kixnw in {6/30/10,7/55/10,8/80/10,9/112/18} {
    \fill[bkVol\kixnv!14, rounded corners=1.5mm] (\kixnx-\kixnw,-3) rectangle (\kixnx+\kixnw,80);
    \node[font=\sffamily\bfseries\footnotesize, text=bkInk] at (\kixnx,77) {第\kixnv 巻};}
  % 路線（横の線）
  \foreach \kixnr/\kixny in {1/68,2/59,3/50,4/41,5/32,6/23,7/14} {
    \draw[bkRoute\kixnr, line width=2.4pt, line cap=round] (8,\kixny) -- (131,\kixny);
    \node[inner sep=0pt] at (1,\kixny) {\bkrouteicon[5.2mm]{\kixnr}};}
  % 合流の印（白い帯。中の黒い点が、入った路線）  x/下/上/記号/点の y
  \foreach \kixnx/\kixnya/\kixnyb/\kixnl/\kixnd in {30/38.5/52.5/ア/{50,41}, 55/29.5/52.5/イ/{50,32}, 80/20.5/70.5/オ/{68,41,23},
      100/38.5/52.5/ア/{50,41}, 108/29.5/52.5/イ/{50,32}, 116/29.5/70.5/ウ/{68,41,32}, 124/11.5/61.5/エ/{59,41,14}} {
    \draw[bkInk, line width=0.9pt, fill=white, rounded corners=2.2mm] (\kixnx-2.6,\kixnya) rectangle (\kixnx+2.6,\kixnyb);
    \node[font=\sffamily\bfseries\footnotesize, text=bkInk] at (\kixnx,\kixnya-3.6) {\kixnl};
    \foreach \kixny in \kixnd {\fill[bkInk] (\kixnx,\kixny) circle (1.1);}}
  \node[anchor=north, kixn lab, text=bkGray, align=center] at (66,-4.5) {（ア）は 6-12 と 9-5・9-9、（イ）は 7-10 と 9-11、（オ）は 8-13・8-14 から 9-0 へ};
\end{tikzpicture}
\bkfigcaption{第6〜9巻での合流の地図。横の線が路線、たての白い帯が合流した所（記号は表\ref{bk:9-13:tab-goryu}）。帯の中の黒い点が、その合流に入った路線（点のない線は、帯の下を通りぬけている）。}
\label{bk:9-13:fig-goryu}
\end{bkfigure}

\bksubsection{5}{（ア）$s=\sigma+j\omega$：倍々と回転が1つの数に}

\label{bk:9-13:s}路線③（倍々）では、$e^{at}$ が「$a$ が正なら倍々にふえ、負なら倍々にへる」ことを学びました（3-11・5-14・6-4）。路線④（回転と波）では、$e^{j\omega t}=\cos\omega t+j\sin\omega t$ が「1秒に $\omega$ ずつ回る矢印」であることを学びました（6-10）。

この2つは、肩の数を複素数 $s=\sigma+j\omega$ にすると、1つの式になります（6-10⑪・9-5）。
\[ e^{st}=e^{(\sigma+j\omega)t}=e^{\sigma t}\,e^{j\omega t}=e^{\sigma t}(\cos\omega t+j\sin\omega t) \]
指数法則（2-7）で2つに分けると、実部 $\sigma$ が「倍々」、虚部 $\omega$ が「回転」を受けもつことがわかります。

\begin{bkfigure}
\kixnmerge{\kixnin{3}{5-14}{$e^{at}$：倍々にふえる・へる}\kixnin{3}{6-4}{冷めるお湯 $20+60e^{-0.05t}$}\kixnin{4}{6-10}{$e^{j\omega t}$：回る矢印}\kixnin{4}{6-12}{$e^{st}$ を入れて特性方程式}}{$e^{st}=e^{\sigma t}(\cos\omega t+j\sin\omega t)$\\[3pt]$\sigma$：倍々、$\omega$：回転}{9-5・9-9}
\bkfigcaption{合流（ア）。$s$ の平面の1つの点（極）が、「どれだけ速く小さくなるか」と「どれだけ速くゆれるか」を同時に教える。}
\end{bkfigure}

6-12 で「同じ $s$」とよんだ特性方程式の解は、9-5 のラプラス変換の $s$ と同じものでした。9-9 では、この点を\textbf{極}とよび、極が $s$ の平面の左半分（$\sigma<0$）にあれば、動きは落ち着くことを学びました。車A の極 $-1,\ -2$ と、車B の極 $-1\pm2j$ は、9-0 の図9で見たとおりです。路線③の終点「制御」は、路線④と合流してはじめて、ゆれる機械にも使える道具になったのです。

\bksubsection{6}{（イ）$e^{At}$：倍々と表が1つの式に}

\label{bk:9-13:eat}路線③では、$\left(1+\dfrac{1}{n}\right)^n$ が、$n$ を大きくすると $e$ に近づくことを見ました（5-14）。路線⑤（表）では、行列をくり返しかけて、少し先の状態を次々に出しました（7-10）。

\begin{bkfigure}
\kixnmerge{\kixnin{3}{5-14}{$\left(1+\frac{x}{n}\right)^n\to e^{x}$}\kixnin{1}{5-17}{$e^{x}=1+x+\frac{x^2}{2!}+\cdots$（5-17 は①の駅）}\kixnin{5}{7-10}{$\left(I+A\frac{t}{n}\right)^n$：行列のくり返し}\kixnin{5}{7-8}{固有値と固有ベクトル}}{$e^{At}=I+At+\dfrac{(At)^2}{2!}+\cdots$\\[3pt]$=\lim\left(I+A\dfrac tn\right)^n$}{9-11}
\bkfigcaption{合流（イ）。数の $e^{x}$ の作り方（5-14・5-17）の $x$ に、行列 $At$ を入れたものが $e^{At}$。}
\end{bkfigure}

冷めるお湯では、7-10 で $\left(1-\dfrac{0.5}{n}\right)^n$ が、$n=10$ で0.5987、$n=1000$ で0.60645 と、$e^{-0.5}=0.60653$ に近づきました。数1つの $e^{-0.5}$ を、行列に広げたものが $e^{At}$ です（9-11）。数1つなら「$a$ が負なら落ち着く」、行列なら「固有値の実部がすべて負なら落ち着く」。倍々の考え（③）が、表の考え（⑤）の中で、もう一度働いています。9-10 の倒立振子で、状態を見て台車を動かすと立ち直るのも、$A-BK$ の固有値を左半分に選んだからでした。

\bksubsection{7}{（ウ）熱の行列の固有ベクトルは、$\sin$ の形}

\label{bk:9-13:sin}8-13 では、長さ \qty{0.1}{m} の金属の棒を9つの点に区切り、熱が広がるようすを、表計算で1歩ずつ計算しました（路線①）。9-1 では、同じ棒の熱を、$\sin$ の波のたし合わせ（フーリエ級数）で解きました（路線④）。そして 9-11 では、8-13 の式を、時間の部分を微分のまま残して $\dot{\vec T}=A\vec T$ と書き、$9\times9$ の行列 $A$ の固有ベクトルを調べました（路線⑤）。

\begin{bkfigure}
\kixnmerge{\kixnin{1}{8-13}{熱方程式を表計算で1歩ずつ}\kixnin{4}{9-1}{$\sin\frac{n\pi x}{L}$ のたし合わせで解く}\kixnin{5}{9-11}{$9\times9$ の行列 $A$ の固有ベクトル}}{いちばん遅い固有ベクトル\\$0.309,\ 0.588,\ 0.809,\ \ldots$\\[2pt]$=\sin\dfrac{m\pi}{10}$（$m=1,\ 2,\ \ldots$）}{9-11}
\bkfigcaption{合流（ウ）。表計算の行列（⑤）の固有ベクトルが、フーリエ級数の $\sin$（④）そのものになる。}
\end{bkfigure}

9-11 で求めた固有ベクトルは、$\sin\dfrac{\pi}{10},\ \sin\dfrac{2\pi}{10},\ \ldots$ と同じ数でした。固有値 $-0.0979$ も、9-1 の1つめの成分が消える速さ $0.0987$〔1/s〕とほとんど同じです。棒のまん中の温度も、10秒後に、9-1 のフーリエ級数で47.45\,℃、9-11 の行列で47.44\,℃と、ほとんど同じ数が出ました。「表で1歩ずつ計算する」ことと「$\sin$ に分けて解く」ことは、同じものを別の方向から見ていたのです。

\bksubsection{8}{（エ）ガウスの形：たし合わせ・はかる・波}

\label{bk:9-13:gauss}つりがねの形 $e^{-x^2}$ は、3つの路線に出てきました。路線⑦（はかる）では、たくさんはかった値のばらつきの形（正規分布：3-17）として。路線②（たし合わせ）では、面積が $\sqrt\pi$ になるガウス積分として（8-6）。そして第9巻では、フーリエ変換しても同じつりがねの形にもどる、いちばん大事な形として（9-4。同じ形にもどる関数は、ほかにもあります）。

\begin{bkfigure}
\kixnmerge{\kixnin{7}{3-14}{標準偏差 $\sigma$：ばらつきの幅}\kixnin{7}{3-17}{正規分布：つりがねの形}\kixnin{2}{8-6}{$\displaystyle\int_{-\infty}^{\infty}e^{-x^2}\,\dd x=\sqrt\pi$}\kixnin{4}{6-10}{回る矢印 $e^{j\omega t}$}\kixnin{7}{9-4}{フーリエ変換：波の成分}}{$e^{-t^2/2}\ \to\ \sqrt{2\pi}\,e^{-\omega^2/2}$\\[3pt]時間の幅が $\sigma$ なら\\周波数の幅は $\dfrac{1}{\sigma}$}{9-4}
\bkfigcaption{合流（エ）。短いパルスほど、周波数は広く広がる。時間の幅と周波数の幅のかけ算は、ほぼ一定。}
\end{bkfigure}

路線⑦の終点は「計測・診断」でした。はかった信号のスペクトルを見て、機械のどこが悪いかを見当づけるとき、「短い時間だけ見ると、周波数ははっきりしない」というこのきまりが、いつもついてまわります。物理シリーズでは、これが電子の位置と運動量（速さ×質量）の「不確定性」（原-3）として、もう一度出てきます。

\bksubsection{9}{（オ）波動方程式と、音・光}

\label{bk:9-13:hado}8-13 では、弦を伝わる波の式（波動方程式）を、場所と時間の偏微分で書きました（路線①）。8-14 では、電気と磁気の4つの式を「わき出し」と「うず」で読み、そこから光が電磁波として速さ $c=1/\sqrt{\mu_0\varepsilon_0}$ で伝わることを見ました（路線⑥）。そして 9-0 で、マイクの所で見た音の波形を、$\sin$ の成分に分けました（路線④）。

\begin{bkfigure}
\kixnmerge{\kixnin{1}{8-13}{波動方程式、答えは $f(x-ct)$}\kixnin{6}{8-14}{4つの式から、光は電磁波}\kixnin{4}{9-0}{1つの場所の波形を $\sin$ に分ける}}{場所と時間で決まる波を\\1つの場所で見ると、\\時間の波形 $p(t)$}{9-0・9-1}
\bkfigcaption{合流（オ）。空間に広がる波（第8巻）を、1つの場所で時間の波形として見ると、第9巻の「分ける」道具が使える。}
\end{bkfigure}

9-1 では、ギターの弦の波動方程式を、$\sin$ の形に分けて解き、倍音が \qty{110}{Hz} の1倍・2倍・3倍…に並ぶわけを見ました。高校の物理で習う弦の固有振動数の式は、波動方程式を分けて解いた答えだったのです。

%% ================= 1つの現象を全巻の道具で =================
\bksection{butsuri}{1つの現象を、全巻の道具で}

\bksubsection{10}{冷めるお湯の長い旅}

\label{bk:9-13:oyu}80\,℃のお湯が、20\,℃の部屋で冷めていく。この1つの場面を、このシリーズでは何回も計算しました。巻が進むたびに、同じお湯を、新しい道具で見直してきたのです。

\begin{bkfigure}
\kixnchain{{2-18}/{表計算で1分ずつ。10分後 55.92\,℃}, {5-14}/{いつも少しずつ冷める先に、$e$ が出てくる}, {6-2}/{オイラー法：今の速さで1歩進む}, {6-4}/{式 $20+60e^{-0.05t}$。\\10分後 56.39\,℃}, {7-10}/{$1\times1$ の行列のくり返し。$0.95^{10}=0.5987$}, {9-5}/{表を逆に引く：\\$\dfrac{20}{s}+\dfrac{60}{s+0.05}$}}{5}
\bkfigcaption{冷めるお湯の旅。路線は ⑤（2-18）→ ③（5-14）→ ①（6-2）→ ③（6-4）→ ⑤（7-10）→ ③（9-5）と、何回も乗りかえている。}
\end{bkfigure}

2-18 の表の55.92\,℃と、6-4 の式の56.39\,℃の差は、「1分の間ずっと、はじめの速さで冷める」としたオイラー法のずれでした。7-10 では、そのずれが $\left(1-\dfrac{0.5}{n}\right)^n$ の $n=10$ と $n\to\infty$ のちがいだとわかり、9-5 では、表を引くだけで、ずれのない答えにたどり着きました。同じお湯でも、道具が変わると、見えるものが変わります。

\bksubsection{11}{車のばねの長い旅}

\label{bk:9-13:kuruma}車のばね（サスペンション）の2つの例、車A（$s^2+3s+2=0$）と車B（$s^2+2s+5=0$）も、長い旅をしました。

どの巻で計算しても、$t=1$ の位置は $2e^{-1}-e^{-2}=0.600$ です。中学の2次方程式の解 $-1$ と $-2$ が、特性方程式の解になり、行列の固有値になり、伝達関数の極になりました。名前は変わっても、同じ2つの数が、車のゆれ方を決めていたのです。

\begin{bkfigure}
\kixnchain{{3-10}/{2次方程式 $s^2+3s+2=0$。解は $-1$ と $-2$}, {6-7}/{ばね・質量・ダンパーの式}, {6-12}/{$x=2e^{-t}-e^{-2t}$（車A）。解は複素平面に}, {7-8}/{状態の行列の固有値は $-1,\ -2$}, {9-7}/{$\dfrac{s+3}{(s+1)(s+2)}$\\を部分分数に}, {9-9}/{極が左半分：落ち着く}, {9-10}/{状態方程式 $\dot{\vec x}=A\vec x$}, {9-11}/{$e^{At}$ で、$t=1$ の位置 0.600}}{7}
\bkfigcaption{車のばねの旅。3-10 で「第6巻では…」と予告した2次方程式が、第9巻の極と状態方程式までつながった。}
\end{bkfigure}

\Needspace{14\baselineskip}\bksubsection{12}{金属の棒の熱と、ロボットの腕}

\label{bk:9-13:bou}金属の棒の熱と、ロボットの腕も、何巻にもまたがって出てきました。

\begin{bkfigure}
\kixnchain{{6-3}/{かべの中の温度。熱は温度の傾きに比例して流れる}, {8-13}/{熱方程式を表計算で1歩ずつ}, {9-1}/{$\sin$ のたし合わせ。まん中が50\,℃は9.47秒}, {9-11}/{熱の行列 $A$ の固有ベクトルは $\sin$ の形}}{3}
\bkfigcaption{金属の棒の熱の旅。6-3 のコラムで紹介したフーリエの考え（1822年）が、9-1 で、フーリエのもとの方法として解けた。}
\end{bkfigure}

\begin{bkfigure}
\kixnchain{{4-8}/{三角関数：回す点の影}, {7-6}/{回転行列で、手の位置 $(0.325,\ 0.391)$}, {8-3}/{ヤコビ行列：関節を少し回すと手はどれだけ動く？}, {9-9}/{腕を動かすときの $\zeta$ を選ぶ}}{3}
\bkfigcaption{ロボットの腕の旅。腕の長さは \qty{0.30}{m} と \qty{0.25}{m}、関節の角は 30°と 45°（7-6・8-3 と同じ）。}
\end{bkfigure}

1つの現象を、1つの巻だけで説明しきることは、ほとんどありません。本物の機械や自然は、いくつもの路線の道具を組み合わせて、はじめて読めるのです。

%% ================= 9-0 の回収 =================
\bksection{goal}{第9巻をふり返る}

\bksubsection{13}{3つの部で、できるようになったこと}

\label{bk:9-13:bu}9-0 では、第9巻を「音を分ける」と「変化を予想する」の2つの柱で始めました。その2つの問いには、もう答えが出ています。「声を成分に分けると、何が見える？」には、9-1〜9-4 で「どの周波数の成分がどれだけあるか（スペクトル）が見え、それが音色の正体」と答えました。「表計算をしないで、10分後や1時間後の温度を予想できないか」には、9-5 で「表を引くだけでできる」と答えました。9-12 では、これらの道具が、音楽や写真の圧縮、ノイズキャンセリング、機械の診断、スマホの通信、MRI、エアコンの PID 制御、ドローンの姿勢の制御、電力の品質（高調波）で働いているようすを読みました（9-12 の「この巻の道具と、働く場面」の表）。%% 9-12（V9-6 の ch9-12.tex に合わせた）

3つの部で、何ができるようになったかをまとめます。

\begin{bkcenter}
\begin{bktable}{colspec={Q[l,m,wd=30mm] Q[l,m,wd=74mm] Q[c,m,wd=18mm]}, rowsep=0.6pt, colsep=3pt, cells={font=\small}, row{1}={font=\sffamily\bfseries\small, bg=bkTheme!10}}
部 & できるようになったこと & 章 \\
第1部\newline フーリエ解析 & くり返す波を $\sin$・$\cos$ に分け（フーリエ級数）、かけて積分して成分を取り出す。くり返さない音も周波数に分け（フーリエ変換）、データは FFT で計算する & 9-1〜9-4 \\
第2部\newline ラプラス変換 & 微分を $s$ のかけ算に翻訳し、表と部分分数で微分方程式を解く。箱のつながりを伝達関数で書き、極の場所で、落ち着くかどうかを見分ける & 9-5〜9-9 \\
第3部\newline 行列と微分方程式 & たくさんの微分方程式を状態方程式にまとめ、$e^{At}$ と固有値で動きを読む。状態を見て入れ方を決める。圧縮・診断・制御での使いみちを読む（9-12） & 9-10〜9-12 \\
\end{bktable}
\bktabcaption{第9巻の3つの部（9-0 の図11と同じ並び）}
\end{bkcenter}

\bksubsection{14}{くり返し出てきた場面の、最後の答え}

\label{bk:9-13:basho}9-0 で「この巻でくり返し出てくる場面」として紹介した主役たちは、それぞれの章で、次のような答えにたどり着きました。

\begin{bkcenter}
\begin{bktable}{colspec={Q[l,m,wd=26mm] Q[l,m,wd=80mm] Q[c,m,wd=16mm]}, rowsep=0.6pt, colsep=3pt, cells={font=\small}, row{1}={font=\sffamily\bfseries\small, bg=bkTheme!10}}
場面 & たどり着いた答え & 章 \\
四角い波 & 4-9 の「約0.785」の正体は $\frac{\pi}{4}$。角のそばには約9\,\% の山が残る（ギブズ現象） & 9-1 \\
ラの音 & FFT で、\qty{440}{Hz}・\qty{880}{Hz}・\qty{1320}{Hz} に 1・0.5・0.25 を見つけた & 9-0・9-4 \\
金属の棒の熱 & まん中が50\,℃になるのは9.47秒。行列の固有ベクトルは $\sin$ の形 & 9-1・9-11 \\
ギターの弦 & A弦の倍音は \qty{110}{Hz} の1倍・2倍・3倍… & 9-1 \\
\end{bktable}
\bktabcaption{第1部の場面の、最後の答え（9-0 の表8）}
\end{bkcenter}

第2部・第3部の場面は、次のとおりです。

\begin{bkcenter}
\begin{bktable}{colspec={Q[l,m,wd=26mm] Q[l,m,wd=80mm] Q[c,m,wd=16mm]}, rowsep=0.6pt, colsep=3pt, cells={font=\small}, row{1}={font=\sffamily\bfseries\small, bg=bkTheme!10}}
場面 & たどり着いた答え & 章 \\
冷めるお湯 & 表を引くだけで $20+60e^{-0.05t}$。10分後 56.39\,℃ & 9-5 \\
RC 回路 & 10\,V で充電すると、2秒で 6.321\,V（$10-10e^{-1}$） & 9-7 \\
車A・車B & 極は $-1,\ -2$ と $-1\pm2j$。どちらも左半分で落ち着く & 9-7・9-9 \\
モーター・エアコン & はかってもどす（フィードバック）と速く追いつくが、比例だけではずれが残る & 9-8・9-9 \\
倒立振子 & 動かさないと2秒で約 0.273\,rad、状態を見て動かすと 0.00092\,rad & 9-10 \\
水そう・送電線 & 水そうは $10\pm10e^{-0.2t}$、送電線の行列から四端子定数 & 9-11 \\
\end{bktable}
\bktabcaption{第2部・第3部の場面の、最後の答え（9-0 の表9）}
\end{bkcenter}

%% ================= 名前と実際の操作 =================
\Needspace{10\baselineskip}\bksection{hon}{名前と、実際の操作の総まとめ}

\bksubsection{15}{全巻の「名前」と「実際の操作」}

\label{bk:9-13:namae}この本を作るときにいちばん大切にしたことの1つは、「記号の名前と、実際にする計算のちがいを、ていねいに見せる」ことでした。第5巻・第8巻・第9巻の記号を、1つの表にまとめます。

\begin{bkcenter}
\begin{bktable}{colspec={Q[c,m,wd=12mm] Q[c,m,wd=28mm] Q[l,m,wd=44mm] Q[l,m,wd=28mm] Q[c,m,wd=11mm]}, rowsep=2pt, colsep=3pt, cells={font=\small}, row{1}={font=\sffamily\bfseries\small, bg=bkTheme!10}}
巻 & 名前（記号） & 名前の意味 & 実際の操作 & 章 \\
第5巻 & $\dfrac{\dd f}{\dd x}$ & $x$ の変わる幅を小さくしていった先の、変わる速さ（傾き） & 導関数の表と、積・合成関数のきまり & 5-4・5-6 \\
第5巻 & $\displaystyle\int_a^b f\,\dd x$ & 細い長方形に分けてたした先 & 原始関数の差 & 5-10・5-11 \\
\end{bktable}
\bktabcaption{第5巻の「名前」と「実際の操作」}
\end{bkcenter}

第8巻では、面や道でたし合わせる積分の名前が出てきました。

\begin{bkcenter}
\begin{bktable}{colspec={Q[c,m,wd=12mm] Q[c,m,wd=28mm] Q[l,m,wd=44mm] Q[l,m,wd=28mm] Q[c,m,wd=11mm]}, rowsep=2pt, colsep=3pt, cells={font=\small}, row{1}={font=\sffamily\bfseries\small, bg=bkTheme!10}}
巻 & 名前（記号） & 名前の意味 & 実際の操作 & 章 \\
第8巻 & $\displaystyle\iint_D f\,\dd A$ & 面をタイルに分けてたした先 & \kixnop{1}（2回） & 8-6 \\
第8巻 & $\displaystyle\int_C\vec F\cdot\dd\vec r$ & 道を小さな1歩に分けてたした先 & \kixnop{1} & 8-7 \\
第8巻 & $\displaystyle\iint_S\vec F\cdot\vec n\,\dd S$ & 向きのついたタイルでたした先 & \kixnop{1}（2回） & 8-8 \\
\end{bktable}
\bktabcaption{第8巻の「名前」と「実際の操作」}
\end{bkcenter}

第5巻では、微分と定積分を、表と原始関数で計算しました。第8巻では、面や道でたし合わせる積分を、すべてふつうの定積分に直しました。第9巻の記号と実際の操作は、9-0 の表6と、巻末の「変換の地図」に、絵つきでまとめてあります。ここでは、記号と道具だけを並べます。

\begin{bkcenter}
\begin{bktable}{colspec={Q[c,m,wd=30mm] Q[l,m,wd=58mm] Q[c,m,wd=22mm]}, rowsep=2pt, colsep=3pt, cells={font=\small}, row{1}={font=\sffamily\bfseries\small, bg=bkTheme!10}}
名前（記号） & 実際の操作 & 章 \\
$a_n$・$b_n$・$c_n$ & \kixnop{1} & 9-1〜9-3 \\
$F(\omega)$ & \kixnop{1}\enspace\kixnop{2} & 9-4 \\
$\mathcal L[f(t)]$ & \kixnop{2} & 9-5・9-6 \\
$\mathcal L^{-1}[F(s)]$ & \kixnop{3}\enspace\kixnop{2} & 9-7 \\
$G(s)$ & \kixnop{2}\enspace\kixnop{4} & 9-8 \\
$e^{At}$ & \kixnop{4}\enspace\kixnop{3} & 9-10・9-11 \\
\end{bktable}
\bktabcaption{第9巻の名前と、実際の操作（くわしくは 9-0 の表6・巻末の「変換の地図」）。$G(s)$ の「行列の計算」は、状態方程式から $C(sI-A)^{-1}B$ で作るとき（表9-8 の E）。}
\end{bkcenter}

2つの表をくらべると、1つのことに気づきます。多くの記号は「細かく分けて、かけて、たした先」か「無限に続くたし算」の名前で、どれも実際は4つの道具のどれかに直して計算しました。第5巻の定積分も、第8巻の面積分も、第9巻のラプラス変換も、名前の意味は「\textbf{細かく分けて、かけて、たした先}」です。そして実際の計算は、どれも「\textbf{表を引く}」か「\textbf{ふつうの定積分に直す}」に行き着きました。第5巻で作った原始関数の表と、第9巻で作ったラプラス変換の表は、同じ役目の道具です。

\[ \kixnNt{\int_0^\infty e^{-0.05t}e^{-st}\,\dd t}=\kixnSt{2}{\frac{1}{s+0.05}} \]

\bksubsection{16}{専門書の式を読む4つの手順}

\label{bk:9-13:tejun}この先、大学の専門書を開くと、見たことのない記号や、見たことのある記号の別の書き方に、たくさん出会います。そのときは、次の4つの手順で読みます。

\begin{enumerate}[label={\bkcircled[\bkTT]{\arabic*}}, leftmargin=2.4\zw, itemsep=3pt, topsep=3pt]
\item \textbf{名前の段で読む}\par\nobreak\noindent 記号の意味を、ことばで言う。「何に・何をかけて・どこからどこまでたすか」「何を何でわるか」。
\item \textbf{記号を訳す}\par\nobreak\noindent 本ごとの書き方のちがいを、この本の書き方に直す（次の⑰の表）。
\item \textbf{実際の操作に直す}\par\nobreak\noindent その章の場合分けの表と流れ図で、4つの道具のどれを使うかを決めて計算する。
\item \textbf{数で確かめる}\par\nobreak\noindent 表計算や Python で、$t=0$ の値や最後の値、いくつかの点の値を計算して、式の答えと同じかを見る。
\end{enumerate}

%% ================= 専門書の読み方ガイド =================
\bksection{card}{専門書の読み方ガイド}

\bksubsection{17}{記号の訳}

\label{bk:9-13:yaku}このシリーズの記号と、専門書でよく見る書き方を、まとめて並べます。\bkidx{せんもんしょのよみかた}{専門書の読み方}同じものが、本によって別の文字や形で書かれていることがあります。逆に、同じ文字が、本によって別のものを表すこともあります。

\begin{bkcenter}
\begin{bktable}{colspec={Q[c,m,wd=26mm] Q[c,m,wd=28mm] Q[l,m,wd=52mm] Q[c,m,wd=14mm]}, rowsep=1.4pt, colsep=3pt, cells={font=\small}, row{1}={font=\sffamily\bfseries\small, bg=bkTheme!10}}
この本 & 専門書の書き方 & 気をつけること & 章 \\
虚数単位 $j$ & $i$ & 数学・物理の本は $i$。電気の本は、電流 $i$ とまぎれないように $j$ & 6-8 \\
$\vec F$ & 太字の $\symbf F$ & 手で書くときは、矢印か、線を1本足した文字 & 4-10・8-14 \\
$\dot x$、$\dfrac{\dd x}{\dd t}$ & $x'$、$\dot x$ & $x'$ は、場所で微分した意味のこともある & 5-1・6-1 \\
\end{bktable}
\bktabcaption{記号の訳（その1：数を表す記号）}
\end{bkcenter}

微分・積分とベクトル場の記号は、次のとおりです。

\begin{bkcenter}
\begin{bktable}{colspec={Q[c,m,wd=26mm] Q[c,m,wd=28mm] Q[l,m,wd=52mm] Q[c,m,wd=14mm]}, rowsep=1.4pt, colsep=3pt, cells={font=\small}, row{1}={font=\sffamily\bfseries\small, bg=bkTheme!10}}
この本 & 専門書の書き方 & 気をつけること & 章 \\
$\dfrac{\partial f}{\partial x}$ & $f_x$、$\partial_xf$ & 下つきの $x$ が「$x$ 成分」の意味の本もある & 8-2 \\
$\mathrm{grad}\,f$、$\mathrm{div}\,\vec F$、$\mathrm{rot}\,\vec F$ & $\nabla f$、$\nabla\cdot\vec F$、$\nabla\times\vec F$、$\mathrm{curl}\,\vec F$ & どれも同じもの。rot と curl も同じ & 8-4・8-9・8-10 \\
$\nabla^2f$ & $\Delta f$ & 差の $\Delta$（$\Delta x$）とは別のもの & 8-12 \\
$\displaystyle\oint_C$、$\displaystyle\oiint_S$ & $\displaystyle\oint_{\partial S}$、$\displaystyle\iint_{\partial V}$ & $\partial S$ は「面 $S$ のふち」、$\partial V$ は「立体 $V$ の表面」 & 8-7・8-8・8-11 \\
固有値 $\lambda$、極 $s$ & $\lambda$、$p$、$s$ & 伝達関数の極は、状態の行列の固有値（打ち消し合う因数がなければ全部） & 7-8・9-9・9-10 \\
\end{bktable}
\bktabcaption{記号の訳（その2：微分・積分・ベクトル場の記号）}
\end{bkcenter}

\begin{bkcenter}
\begin{bktable}{colspec={Q[c,m,wd=26mm] Q[c,m,wd=28mm] Q[l,m,wd=52mm] Q[c,m,wd=14mm]}, rowsep=1.4pt, colsep=3pt, cells={font=\small}, row{1}={font=\sffamily\bfseries\small, bg=bkTheme!10}}
この本 & 専門書の書き方 & 気をつけること & 章 \\
$\dfrac{a_0}{2}$（平均） & $a_0$ & $a_0$ を平均そのものとする本もある。係数の式の前の数（$\frac1T$ か $\frac2T$）を確かめる & 9-1・9-2 \\
もどす式に $\dfrac{1}{2\pi}$ & 両方の式に $\dfrac{1}{\sqrt{2\pi}}$ & $2\pi$ の置き場所が本でちがう。周波数（Hz）で書いて $2\pi$ を肩に入れる本もある & 9-4 \\
$\mathcal L[f(t)]$ & $\mathcal L\{f(t)\}$ & かっこの形がちがうだけ & 9-5 \\
$\mathcal L^{-1}$ & 複素数の積分の式 & 名前の上の式。実際は部分分数と表 & 9-7 \\
たたみこみの $t_1$ & $\tau$ & この本の $\tau$ は時定数 & 9-6 \\
時定数 $\tau$ & $T$ & $T$ は周期・温度にも使う文字 & 6-4・9-8 \\
$e^{At}$ & $\exp(At)$、$\Phi(t)$ & 同じもの。$\Phi(t)$ は「状態遷移行列」とよぶ本も & 9-11 \\
電場の流束 $\Phi_E$ & 電束 $\Psi$（別の量） & 真空中では $\Psi=\varepsilon_0\Phi_E$。電気工学の「電束」は $\vec D$ の流束（単位 C） & 8-8・8-14 \\
$\mu_0\approx4\pi\times10^{-7}$ N/A² & $\mu_0=4\pi\times10^{-7}$ N/A² & 2019年の SI の改定からは、測って決める値で、ほぼこの値 & 8-10・8-14 \\
$\delta(t)$ & $\delta(t)$ & 6-7 の対数減衰率 $\delta$ とは別のもの & 9-6 \\
\end{bktable}
\bktabcaption{記号の訳（その3：フーリエ・ラプラス・電磁気の記号）。信号処理・制御の本の記号（$X(\omega)$、$H(s)$、$u(t)$、ずれの $e(t)$ など）は、9-12 の「専門書をのぞいてみよう」の表を見る。}
\end{bkcenter}

\bksubsection{18}{専門書の式を1つ読んでみよう}

\label{bk:9-13:rei}電気回路の本には、たとえば次のような式がのっています。抵抗 $R$ とコンデンサー $C$ を直列につなぎ、スイッチを入れて電圧 $E$ をかけたときの式です。
\[ \frac{E}{s}=\Bigl(R+\frac{1}{sC}\Bigr)I(s) \]
これを、⑯の4つの手順で読んでみます。数は 6-4・9-7 と同じ $E=\qty{10}{V}$、$R=\qty{20}{k\ohm}$、$C=\qty{100}{\micro F}$ にします。

\begin{enumerate}[label={\bkcircled[\bkTT]{\arabic*}}, leftmargin=2.4\zw, itemsep=3pt, topsep=3pt]
\item \textbf{名前の段}\par\nobreak\noindent 大文字の $I(s)$ は、電流 $i(t)$ のラプラス変換（名前）。$\dfrac{E}{s}$ は、スイッチを入れてからずっと $E$ の電圧（階段の形）の変換。$\dfrac{1}{sC}I(s)$ は、コンデンサーの電圧 $\dfrac{1}{C}\displaystyle\int_0^t i\,\dd t_1$ の変換（積分は $s$ でわる：9-6）。
\item \textbf{記号を訳す}\par\nobreak\noindent この本と同じ書き方。はじめにコンデンサーに電気がない（はじめの値が0）ことが、式に入っている。
\item \textbf{実際の操作}\par\nobreak\noindent かけ算・わり算で $I(s)$ を求め、表を引く。
\begin{bkhenkei}[0.5]
I(s)=\dfrac{E}{s}\div\Bigl(R+\dfrac{1}{sC}\Bigr) & 両辺を $R+\frac{1}{sC}$ でわる \\
=\dfrac{E}{sR+\frac{1}{C}} & $\frac{E}{s}$ と $R+\frac{1}{sC}$ の両方に $s$ をかける \\
=\dfrac{E}{R}\cdot\dfrac{1}{s+\frac{1}{RC}} & 分母を $R$ でくくる \\
i(t)=\dfrac{E}{R}\,e^{-t/(RC)} & 表の $\frac{1}{s+a}\to e^{-at}$（9-5） \\
\end{bkhenkei}
\item \textbf{数で確かめる}\par\nobreak\noindent 電流のはじめの値と、時定数は
\[ \frac{E}{R}=\frac{\qty{10}{V}}{\qty{20000}{\ohm}}=\qty{0.5}{mA},\qquad RC=\qty{20000}{\ohm}\times\qty{0.0001}{F}=\qty{2}{s} \]
です。$t=0$ で \qty{0.5}{mA}、2秒後は $0.5\times e^{-1}=\qty{0.184}{mA}$ です。抵抗の電圧は $Ri=10e^{-t/2}$ なので、コンデンサーの電圧は $10-10e^{-t/2}$ です。2秒で \qty{6.321}{V} となり、9-7 の答えと同じです。
\end{enumerate}

専門書の式も、1行ずつ「名前」と「実際の操作」に分けて読めば、この本で学んだ道具だけで読めます。

\bksubsection{19}{大学の科目と、この本の巻}

\label{bk:9-13:kamoku}大学の工学部では、数学や、数学を使う専門の科目を、たくさん学びます。科目の名前と、このシリーズのどの巻・章がその土台になっているかを並べます。専門書を開いて手が止まったら、右の列の章にもどってみましょう。

\begin{bkcenter}
\begin{bktable}{colspec={Q[l,m,wd=40mm] Q[l,m,wd=50mm] Q[l,m,wd=32mm]}, rowsep=0.5pt, colsep=3pt, cells={font=\small}, row{1}={font=\sffamily\bfseries\small, bg=bkTheme!10}}
大学の科目 & 中心になる考え & この本の巻・章 \\
微分積分学 & 微分・積分・偏微分・重積分 & 第5巻、8-1〜8-6 \\
線形代数 & 行列・固有値・対角化 & 第7巻 \\
微分方程式 & 特性方程式、連立の式 & 第6巻 \\
ベクトル解析 & 勾配・発散・回転、2つの定理 & 第8巻 \\
応用数学 & フーリエ変換・ラプラス変換 & 第9巻 \\
確率・統計 & ばらつき、最小二乗法、正規分布 & 3-13〜3-17 \\
数値計算 & 差分、オイラー法 & 2-18・6-2・7-10・8-13 \\
電気回路 & $e^{-t/\tau}$、交流の複素数、伝達関数 & 6-4・6-10・9-7・9-8 \\
制御工学 & 伝達関数・極・状態方程式 & 9-5〜9-11 \\
信号処理 & スペクトル・FFT・サンプリング & 9-1〜9-4・9-12 \\
機械力学・振動工学 & ばねとダンパー、共振 & 6-6・6-7・6-12・6-13 \\
流体力学・伝熱工学 & 流束・発散、熱方程式 & 8-8・8-9・8-12・8-13 \\
電磁気学 & 2つの定理、4つの式 & 8-9〜8-11・8-14 \\
\end{bktable}
\bktabcaption{大学の科目と、その土台になる巻・章（科目の名前や中身は、大学によって少しずつちがう）}
\end{bkcenter}

どの科目も、1つの巻だけでなく、いくつもの路線の道具を組み合わせて使います。目次や索引で見知らぬことばに出会ったら、まずこのシリーズのさくいんで、同じことばをさがしてみてください。

%% ================= 物理シリーズへの橋 =================
\bksection{butsuri}{物理シリーズへの橋}

\bksubsection{20}{数学の道具を、物理で使う}

\label{bk:9-13:butsuri}このシリーズの数学は、「物理や工学を説明するための数学」でした。この先は、姉妹編の物理シリーズ（全5巻）で、手に入れた道具を使って、機械・熱・波・電気・原子を読んでいきます。\bkidx{ぶつりしりーず}{物理シリーズ}物理の本の章には、その章で使う数学の章の番号が書いてあります。

\begin{bkcenter}
\begin{bktable}{colspec={Q[c,m,wd=14mm] Q[l,m,wd=20mm] Q[l,m,wd=40mm] Q[l,m,wd=46mm]}, rowsep=1.4pt, colsep=3pt, cells={font=\small}, row{1}={font=\sffamily\bfseries\small, bg=bkTheme!10}}
巻 & 題名 & 副題 & とくによく使う数学 \\
第1巻 & 力学 & 動く・回る・ゆがむ・流れる & 第8巻（ベクトル解析）、第7巻（行列） \\
第2巻 & 熱学 & エネルギーを使いこなす & 第8巻、第3巻・第5巻 \\
第3巻 & 波動 & ゆれと波を読む & \textbf{第9巻}、第6巻（微分方程式） \\
第4巻 & 電磁気学 & 場から回路・機械・電波へ & 第8巻、第6巻 \\
第5巻 & 原子 & 量子から半導体・原子力へ & 第6巻、第4巻ほか \\
\end{bktable}
\bktabcaption{物理シリーズ全5巻と、とくによく使う数学の巻（それぞれの章の「前に学ぶ数学の章」を数えたもの）}
\end{bkcenter}

第9巻の道具は、とくに物理シリーズ第3巻（波動）で、たくさん使います。

\begin{bkcenter}
\begin{bktable}{colspec={Q[c,m,wd=12mm] Q[l,m,wd=60mm] Q[l,m,wd=46mm]}, rowsep=1.4pt, colsep=3pt, cells={font=\small}, row{1}={font=\sffamily\bfseries\small, bg=bkTheme!10}}
物理の章 & 題名 & 使う第9巻の章 \\
波-2 & ゆすり方で変わるゆれ：強制振動と共振 & 9-8（周波数応答・dB） \\
波-4 & ねらいどおりに動かす：フィードバック制御 & 9-6・9-8・9-9・9-10 \\
波-6 & はね返って重なる：反射・定在波・固有振動 & 9-1・9-2・9-3（変数分離） \\
波-10 & 見えない中を調べる：超音波と振動の計測 & 9-4（スペクトル） \\
波-11 & 音を数字にする：信号とスペクトル & 9-4（FFT・サンプリング） \\
波-13 & 光の波：干渉・回折・偏光・散乱 & 9-4 \\
波-14 & 波で情報を送る & 9-12 \\
熱-11 & 壁を通りぬける熱：熱伝導と熱抵抗 & 9-1（変数分離） \\
電-9 & スイッチを入れた瞬間：過渡現象 & 9-7 \\
原-3 & 電子は波でもある：ド・ブロイ波と不確定性 & 9-4（ガウスの形） \\
原-5 & 波で書く電子：シュレディンガー方程式 & 9-1・9-3 \\
原-15 & 原子炉の中の中性子：核分裂と原子炉 & 9-9（極と安定） \\
\end{bktable}
\bktabcaption{物理シリーズで、第9巻の道具を使う章}
\end{bkcenter}

物理シリーズでは、「こまったことがあれば、この数学の章にもどろう」と、章の番号で数学の本を指しています。物理の本を読みながら、この9冊を、いつでも開ける道具箱として使ってください。

\begin{furikaeri}
\begin{checklist}
\item 7本の路線の名前と終点を言い、全9巻の路線図でたどれる。\bkfurref{bk:9-13:rosen}
\item 路線が合流した5つの所を、式と例で説明できる。\bkfurref{bk:9-13:goryu}
\item 専門書の式を、名前の段で読み、記号を訳し、4つの道具で計算できる。\bkfurref{bk:9-13:tejun}
\end{checklist}
\end{furikaeri}

\begin{column}{数学は、何のために学ぶのか}
このシリーズの1巻めのはじめに、「かずを つかうと、なにが いいの？」（1-0）と問いかけました。「たくさん」より「12こ」のほうが、よく伝わる。数を使うと、ものごとを、だれにでも同じように伝えられる。それが、最初の答えでした。

9冊を読んだいま、もう少しくわしく答えられます。数学は、自然や機械のふるまいを、短いことばで書き、まだ起きていないことを予想するための道具です。音を成分に分けて、どこが悪いかを見つける。お湯や車や電気回路が、この先どう動くかを、表を引いて予想する。ロボットやドローンが、たおれずに落ち着くように、固有値を選ぶ。どれも、小学1年の「数える」から一歩ずつ積み上げた道具で、できるようになったことです。
\end{column}

\begin{bkfigure}
\begin{tikzpicture}[x=1mm, y=1mm]
  \kixntalk{0}{sora}{ソラ}{9冊、全部読み終わったね。最初は、リンゴを数えるところからだったのに。}
  \kixntalk{-19}{mina}{ミナ}{同じ冷めるお湯を、表計算でも、$e$ でも、行列でも、ラプラス変換でも計算したね。道具がふえるたびに、見えるものがふえていったよ。}
  \kixntalk{-41}{delta}{デルタ}{ここまで、本当によくがんばったね。全部を一度で覚えていなくても大丈夫。困ったときは、章の番号をたよりに、いつでももどってこられるよ。この道具箱を持って、物理と工学の世界へ出かけよう。}
\end{tikzpicture}
\end{bkfigure}

これで、『物理を理解するための数学』全9巻はおしまいです。ここまで読み通したあなたは、もう、工学の専門書を自分の力で読み進められます。わからない所に出会ったら、この地図を開いて、どの路線の、どの駅の道具かを確かめてください。

\begin{bkfigure}
\begin{tikzpicture}[x=1mm, y=1mm]
  \node[font=\sffamily\bfseries\small, text=bkInk, anchor=west] at (0,136) {物理を理解するための数学（全9巻）};
  \foreach \kixnv/\kixnt in {1/かぞえる・はかる,2/数と量のことば,3/関数とデータ,4/形と向き,5/変化をつかむ,6/未来を計算する,7/まとめて計算する,8/空間に広がる量,9/波と変化を読みとく} {
    \pgfmathtruncatemacro\kixnr{(\kixnv-1)/3}\pgfmathtruncatemacro\kixnk{mod(\kixnv-1,3)}
    \fill[bkVol\kixnv, rounded corners=1mm] ({\kixnk*44.5},{130-\kixnr*20}) rectangle ++(9,-17);
    \node[text=bkBandOn\kixnv, font=\sffamily\bfseries\normalsize] at ({\kixnk*44.5+4.5},{121.5-\kixnr*20}) {\kixnv};
    \draw[bkVol\kixnv, line width=0.9pt, rounded corners=1mm] ({\kixnk*44.5},{130-\kixnr*20}) rectangle ++(42,-17);
    \node[anchor=west, font=\sffamily\bfseries\footnotesize, text=bkInk] at ({\kixnk*44.5+11},{121.5-\kixnr*20}) {\kixnt};}
  \draw[kixn arr, line width=1.6pt] (64,66) -- (64,28);
  \bkChar{sora}{(32,30)}{30mm}
  \bkChar{mina}{(49,30)}{30mm}
  \bkChar{delta}{(88,30)}{25mm}
  \foreach \kixnt/\kixns [count=\kixni from 0] in {力学/{動く・回る・\\ゆがむ・流れる},熱学/{エネルギーを\\使いこなす},波動/{ゆれと波を\\読む},電磁気学/{場から回路・\\機械・電波へ},原子/{量子から\\半導体・\\原子力へ}} {
    \node[draw=bkInk!60, line width=0.9pt, rounded corners=1.5mm, fill=white, minimum width=24mm, minimum height=10mm,
      font=\sffamily\bfseries\small, text=bkInk] (kixnp\kixni) at ({\kixni*26.5+12.5},18) {\kixnt};
    \node[anchor=north, text width=24mm, align=flush center, font=\footnotesize, text=bkInk, inner sep=0pt] at ([yshift=-1.8mm]kixnp\kixni.south) {\kixns};}
  \node[font=\sffamily\bfseries\small, text=bkInk, anchor=west] at (0,-10) {物理シリーズ（全5巻）};
\end{tikzpicture}
\bkfigcaption{数学の9冊から、物理の5冊へ。ソラ・ミナ・デルタといっしょに、次の本へ。}
\end{bkfigure}
